\documentclass{article}
\usepackage{microtype}
\usepackage{siunitx}
\usepackage{amsmath}

\title{art class questions}
\author{ahaan shokeen}
\date{2026-04-21}

\begin{document}
  \maketitle

  \section{strand \( \emptyset \)}

  \begin{enumerate}

    \item the 1960s were a time of great global transformation. the cold war created tensions between the united states of america and the union of soviet socialist republics. artists realized that normal, traditional art could not accurately represent these times of turmoil. they decided that art was not about what it looked like, but about the meaning. they asked the question: \emph{what does this mean}? and questioned things like who decided the definition of art.
    \item because the focus of the art was now the idea and meaning, the technical skill behind making the artwork did not matter as much. now, artists did not simply make a nice picture, they expressed  their thoughts, emotions, and ideas through subtext or an implied meaning. people were now expected to analyse, question and think about the art they saw, rather than just look at it and admire the technical skills.
    \item in conceptualism, changing the context has a massive effect. because conceptual art is \emph{heavily symbolic}, the way it is represented changes the artwork entirely. for example, the artwork \emph{one and three chairs} by joseph kosuth shows how the same word, \emph{chair}, can be interpreted differently as a picture of a chair, a real chair, or the definition of a chair.


  \end{enumerate}

\end{document}

